\section{Introduction}\label{sec:introduction}

The abundance problem asks when the numerical positivity of a log canonical
bundle produces a holomorphic map. More precisely, a nef adjoint should have
a positive multiple generated by global sections. On a compact Kähler space,
nefness is an analytic condition on a cohomology class, whereas generation
is a statement about an actual holomorphic line bundle. The distinction is
essential in the nonprojective setting.

We prove log abundance for compact Kähler spaces under logarithmic Iitaka
subadditivity for smooth projective varieties with reduced boundary. The
argument passes through a stronger birational statement: every
pseudo-effective smooth log adjoint has a semiample positive part and an
effective rational fixed divisor equal to its analytic divisorial negative
part. This form both supports the induction and retains the line-bundle
information needed for generation on the original space.

\subsection{The assumption and the theorem}

A compact Kähler space is a compact complex analytic space with a Kähler
form in the sense of local strictly plurisubharmonic potentials on local
embeddings. For a rational line bundle \(L\), analytic nefness means that
\(c_1(L)\) belongs to the closure of the Kähler cone. Equivalently, fix a
Kähler form \(\omega\) and an integer \(r>0\) for which \(rL\) is a line
bundle. For every \(\varepsilon>0\), that line has a smooth Hermitian metric
\(h_\varepsilon\) satisfying
\[
 \frac{\sqrt{-1}}{2\pi r}\Theta_{h_\varepsilon}
 \geq-\varepsilon\omega.
\]
Smooth metrics and forms on a singular space are understood through local
embeddings. We call \(L\) semiample if an actual positive integral multiple
is a holomorphic line bundle generated by its global sections.

\begin{assumption}[Logarithmic Iitaka subadditivity]
\label{asm:log-iitaka}
Let \(f:X\to Y\) be a surjective morphism with connected fibers between
smooth connected projective complex varieties. Let \(D_X\) and \(D_Y\) be
reduced effective simple normal crossing divisors, allowing zero divisors,
such that
\[
 \Supp(f^*D_Y)\subseteq\Supp(D_X).
\]
For a very general smooth fiber \(F\), put \(D_F=D_X|_F\). Then
\[
 \kappa(X,K_X+D_X)\geq
 \kappa(F,K_F+D_F)+\kappa(Y,K_Y+D_Y).
\]
Here \(D_F\) is reduced SNC and
\(K_F+D_F\sim (K_X+D_X)|_F\). The Iitaka dimension is \(-\infty\) when
all positive integral systems are empty, with
\((-\infty)+b=-\infty\) also for \(b=-\infty\); the zero divisor on a
point has Iitaka dimension zero.
\end{assumption}

\begin{theorem}[Log abundance for compact Kähler spaces]\label{thm:main}
Assume Assumption~\ref{asm:log-iitaka}. Let \(X\) be a normal irreducible
compact Kähler complex analytic space, and let \(\Delta\geq0\) be a rational
divisor such that \((X,\Delta)\) is log canonical and \(K_X+\Delta\) is
rational Cartier. If \(K_X+\Delta\) is analytically nef, then it is
semiample. Explicitly, there is an integer \(m>0\) such that
\(m(K_X+\Delta)\) is Cartier and
\[
 H^0\bigl(X,\OO_X(m(K_X+\Delta))\bigr)\otimes_{\C}\OO_X
 \longrightarrow \OO_X(m(K_X+\Delta))
\]
is surjective at every point of \(X\).
\end{theorem}

The theorem includes every finite dimension and the zero-boundary case.
Assumption~\ref{asm:log-iitaka} is used through the projective good-model
theorem stated in Proposition~\ref{proj:good-models}. The cited proof of
that projective theorem uses only its zero-boundary case, on a resolved
projective Albanese fibration.

\subsection{Context and relation to earlier work}

In the minimal model program, abundance complements the construction of
minimal models: the nef adjoint on a minimal model should determine a
canonical fibration. For minimal projective threefolds, Miyaoka proved the
case of numerical dimension one and Kawamata completed canonical abundance
\cite{Miyaoka1988,Kawamata1992}. Keel, Matsuki, and McKernan established
log abundance for threefolds
\cite{KMM1994,KMM2004}. In arbitrary dimension, Birkar, Cascini, Hacon, and
McKernan established log terminal models for projective klt pairs with big
boundary and pseudo-effective adjoint, and finite generation for big
adjoints \cite{BCHM2010}.

The compact Kähler history already points to the importance of simple
spaces. Peternell's threefold theorem isolated the possible case of a simple
space not bimeromorphic to a finite quotient of a torus
\cite{Peternell2001}; Demailly and Peternell subsequently proved canonical
nonvanishing for nef terminal Kähler threefolds, including that case
\cite{DemaillyPeternell2003}. Höring and Peternell constructed minimal models
for normal \(\Q\)-factorial terminal compact Kähler threefolds with pseudo-effective canonical class
\cite{HoeringPeternell2016}. Campana, Höring, and Peternell treated canonical abundance
\cite{CHP2016,CHP2023}, and Das and Ou
established log abundance for compact Kähler log canonical threefolds
\cite[Corollary~1.3]{DasOu2026}.

With \(D_X=D_Y=0\), Assumption~\ref{asm:log-iitaka} is Iitaka's classical
\(C_{n,m}\) inequality
\(\kappa(X)\geq\kappa(F)+\kappa(Y)\) for algebraic fiber spaces
\cite[Equation~(1.0.1)]{CaoPaun2017}. Cao and Păun proved the analogous
logarithmic inequality over an abelian base when the pair on the total space
has an effective rational boundary and is klt
\cite[Theorem~1.1]{CaoPaun2017}. Fujino derives
logarithmic subadditivity from the conjectured equality
\(\kappa_\sigma(X,K_X+D)=\kappa(X,K_X+D)\) for smooth projective
\(X\) with reduced SNC \(D\)
\cite{Fujino2017Subadditivity,Fujino2020Corrigendum}. Hashizume proves the
reduced-SNC subadditivity statement when the log canonical divisor of the
very general fiber is abundant \cite[Theorem~1.2]{Hashizume2020}. Here the
implication is used in the other direction: the projective companion
\cite{LI} derives good models for pseudo-effective projective lc adjoints
from Assumption~\ref{asm:log-iitaka}, and those good models anchor the
Kähler induction.

Hacon and Xie's cone theorem, adjunction, relative positivity, and
modified-big preparation \cite{HX} provide analytic inputs for our
program constructions. Section~\ref{sec:programs} constructs effective
generalized base data on a higher carrier, with an exact adjoint
pullback identity and a multiplier-ideal inclusion. Fibrewise
$L^2$ positivity and regularization give this presentation; its
Stein-factor consequence supplies the non-klt induction.
The section then proves the restricted semiampleness, proper relative
good-model, and finite-model statements needed here, following the
dimension-induction strategy of \cite{HX}.
The nef data remain globally nef on a fixed carrier, and the total
generalized boundary is globally modified big. Actual rational lines
make the selected ray contractions projective; ordinary projective
analytic results \cite{FujinoAnalyticBCHM} then construct their flips
and the relative models. This produces a Kähler pullback description of
the relevant Bott--Chern class. Generation of the prescribed holomorphic
line requires the additional divisorial decomposition and boundary
arguments. The fibration comparison and period positivity come from
\cite{CI}; the boundary and lifting constructions adapt inputs in
\cite{CB,BL}. The precise imported results are stated at their points
of use.

\subsection{The proof and its main constructions}

For a pseudo-effective \((1,1)\)-class \(\alpha\), \(N(\alpha)\) records
the least generic vanishing forced along prime divisors. Given a smooth
compact Kähler manifold \(X\) and a rational SNC boundary \(B\) with
\(J=K_X+B\) pseudo-effective, the induction constructs a modification
\(\mu:Y\to X\) with \(Y\) smooth and
\[
 \mu^*J\sim_{\Q}P+R,\qquad P\text{ semiample},\qquad
 R=N(c_1(\mu^*J)),
\]
as an identity of rational holomorphic line bundles, with \(R\) an effective
rational divisor. Section~\ref{sec:inductive-reduction} uses transfer rules,
algebraic reduction, and generically finite covers to reduce the induction
to projective manifolds, nontrivial fibrations, and simple spaces of
algebraic dimension zero. Here simple means that no positive-dimensional
proper compact subvariety passes through a very general point.

Section~\ref{sec:programs} supplies two program results. For the fibration
case it contracts a known negative part while preserving the prescribed nef
line. For the simple case it proves special termination for a chosen scaling
and reaches a nef model when the pseudo-effective adjoint is rationally
equivalent to a divisor supported on the reduced boundary, allowing negative
coefficients. The termination argument uses lower-dimensional induction to
construct one model on which an interval of perturbed adjoints is nef. On a
common resolution their negative multiplicities are affine, whereas every
nontrivial wall of the restricted program changes a slope. Hence only
finitely many such walls occur.

Section~\ref{sec:boundary} turns generation on separate strata into generation
on the entire reduced dlt boundary, adapting Fujino's method of admissible
sections \cite[Section~4]{Fujino2000}. When a comparison is needed, a Mori
contraction realizes Kollár's residue comparison at the two coefficient-one
points of a general \(\PP^1\) fiber
\cite[Definition~13 and Proposition~14]{Kollar2013Sources}. On a common
resolution, the pulled-back restrictions of the ambient Kähler class differ
by a real linear combination of Chern classes of line bundles. On a stratum
whose boundary dominates the image of the map
defined by its semiample adjoint, restriction determines sections. For the
remaining strata, we prove finite image for self-comparisons on sufficiently
divisible pluricanonical systems using an invariant integral, period and
lattice arguments, and a uniform cohomological bound for cyclic covers.
Products over these finite images give compatible generating sections, used
in both closing arguments.

Section~\ref{sec:rank-one} first reduces to fibrations whose very general
fiber has log Kodaira dimension zero. On a prepared fibration \(g:X\to W\),
a relative generating form gives
\[
 K_X+B\sim_{\Q}g^*H+A_*,\qquad H=K_W+T+M,
\]
where \(T\) is a rational SNC boundary and \(M\) is pulled back from a nef
rational line on a projective quotient of \(W\). Fiber induction allows us
to subtract the part of \(A_*\) dominating \(W\) from a positive current for
\(K_X+B\). The resulting metric descends along connected smooth fibers; a
zero among the nonnegative coefficients at primes dominating each base prime
allows extension, proving \(H\) pseudo-effective. If \(a(W)=0\), the
projective quotient is a point, so \(M\sim_{\Q}0\) and lower-dimensional
induction makes \(H\) rationally equivalent to its negative divisor. For
projective \(W\), a reduction via a chosen generalized program derived from
Assumption~\ref{asm:log-iitaka} either lowers the positive base dimension,
handled by secondary induction, or reaches a nef line that is big or torsion.
At a stopping case, an intersection argument identifies the full negative
divisor upstairs. A torsion positive part completes the decomposition
directly. In the big nef case, Section~\ref{sec:programs} contracts that
divisor while preserving the nef line; boundary generation, extension, and
a new log canonical place at a hypothetical base locus prove semiampleness.

Section~\ref{sec:nonvanishing} proves, independently of
Assumption~\ref{asm:log-iitaka}, that a positive canonical power has a nonzero
meromorphic section on every smooth connected simple compact Kähler manifold
of algebraic dimension zero. Ou's uniruledness and foliation results supply
the cotangent slope control \cite{Ou2025}. A point-threshold bound for big
classes of volume one is contradicted using an auxiliary projective bundle
over \(X\times X\). The diagonal in its square produces a subsheaf occupying
a fixed positive fraction of a symmetric cotangent power of that bundle.
A second construction over the diagonal of \(X\times X\) forces determinant
vanishing which, after restriction to a blowup of \(X\), violates the point
bound.

Finally, Section~\ref{sec:signed} uses the meromorphic section to obtain,
after resolving and enlarging to a reduced SNC boundary, a pseudo-effective
adjoint with a signed representative supported there.
Section~\ref{sec:programs} gives a nef dlt model, and
Section~\ref{sec:boundary} gives generation on its reduced boundary.
Pseudo-effectivity of the canonical pullback and an intersection argument
isolate positive-dimensional boundary fibers away from components whose
signed coefficients are nonpositive. A local root construction adapted from
\cite[Proposition~3.3]{LI} separates the positive and negative divisor
supports. An SNC Hodge-module calculation extends the lifting method of
\cite[Sections~10--12]{BL} to residual poles and kills every finite-order
obstruction to lifting such a fiber. Douady space and Artin approximation
give compact deformations leaving the boundary; compactness of the cycle-space
components and relative compactness of bounded-volume cycles allow a Baire
argument to produce a covering family. Simplicity
forces the actual nef line to be torsion. The program comparison and
Lemma~\ref{bir:unique-current} then permit subtraction of the added boundary,
yielding the inductive decomposition.
